% !PRINTED 112
\weberpage{112}
%<<CONT

eigentlichen Algebra fremd und darum ist es befriedigender, die \gesperrt{principiellen} Fragen, so wie es im Vorhergehenden geschehen ist und auch noch später geschehen soll, ohne ihre Hülfe zu beantworten. Für die Anwendung und eine bequemere Anschauung wollen wir aber dieses Hülfsmittel doch nicht entbehren.

Setzen wir, wenn $b,c$ reelle Zahlen sind:
\begin{equation}
 b=r\cos\varphi,\qquad c=r\sin\varphi,
\tag{1}
\end{equation}
so ist, wenn wir $r$ positiv annehmen, der Winkel $\varphi$ hierdurch bis auf ein Vielfaches von $2\pi$ bestimmt. Er wird völlig bestimmt sein, wenn wir ein Intervall von der Grösse $2\pi$ festsetzen, in dem er liegen soll, z. B.:
\[
 -\pi<\varphi\le \pi.
\]
In geometrischer Auffassung sind $r,\varphi$ die Polarcoordinaten des Punktes, dessen rechtwinklige Coordinaten $b,c$ sind. $r$ ist der absolute Werth der complexen Grösse $a=b+ci$ und $\varphi$ wollen wir die \gesperrt{Phase} dieser complexen Grösse nennen.

Das Product zweier complexen Grössen
\begin{equation}
\begin{aligned}
 a&=r(\cos\varphi+i\sin\varphi),\\
 a'&=r'(\cos\varphi'+i\sin\varphi'),
\end{aligned}
\tag{2}
\end{equation}
ergiebt sich in der Form
\begin{equation}
 aa'=rr'[\cos(\varphi+\varphi')+i\sin(\varphi+\varphi')],
\tag{3}
\end{equation}
und daraus erhält man, wenn $n$ eine beliebige ganze positive (oder selbst negative) Zahl ist, den nach \gesperrt{Moivre} benannten Lehrsatz:
\begin{equation}
 (\cos\varphi+i\sin\varphi)^n=\cos n\varphi+i\sin n\varphi,
\tag{4}
\end{equation}
der dazu geführt hat, die Grösse $\cos\varphi+i\sin\varphi$ als Exponentialfunction mit imaginärem Exponenten zu betrachten, und demnach
\[
 \cos\varphi+i\sin\varphi=e^{i\varphi}
\]
zu setzen. Dadurch erhalten die Formeln (3) und (4) die Gestalt
\begin{equation}
 e^{i\varphi}e^{i\varphi'}=e^{i(\varphi+\varphi')},\qquad (e^{i\varphi})^n=e^{in\varphi}.
\tag{5}
\end{equation}
Wenn wir demnach
\begin{equation}
 a=r(\cos\varphi+i\sin\varphi)
\tag{6}
\end{equation}
setzen und unter $\sqrt[n]{r}$ den nach §. 32, IX. existirenden und völlig bestimmten positiven Werth verstehen, dessen $n$te Potenz $=r$ ist, so ist
%CONT>>
